% =====================================================================
%  main.tex -- AMS-style manuscript.
%  Note the AMS ordering: abstract goes BEFORE \maketitle, and
%  \subjclass / \keywords come after \maketitle.
%  Build with:  pdflatex main; bibtex main; pdflatex main; pdflatex main
% =====================================================================
\documentclass[11pt,reqno]{amsart}
% =====================================================================
%  preamble.tex -- AMS-style variant.
%  amsart already supplies theorem environments (thm, lem, prop, cor,
%  defn, rem, conj, ...) and handles \title / \author / \date / abstract
%  itself, so this preamble is much shorter than the article one.
% =====================================================================

\usepackage[T1]{fontenc}
\usepackage[utf8]{inputenc}
\usepackage{lmodern}
\usepackage{amsmath,amssymb,amsthm,mathrsfs}
\usepackage{geometry}
\geometry{a4paper,margin=30mm}
\usepackage{microtype}
\usepackage{graphicx,booktabs,array}
\usepackage{enumitem}
\setlist{itemsep=3pt,topsep=5pt}
\usepackage{xcolor}
\definecolor{linkblue}{RGB}{31,70,105}
\usepackage{tikz}
\usetikzlibrary{arrows.meta,positioning,calc,decorations.pathmorphing,fit}

% hyperref must come after everything else that defines commands.
\usepackage[colorlinks=true,linkcolor=linkblue,citecolor=linkblue,
            urlcolor=linkblue]{hyperref}
\usepackage[capitalise,noabbrev,nameinlink]{cleveref}

\numberwithin{equation}{section}

% amsart numbers theorem-like environments continuously (1, 2, 3, ...).
% Uncomment the next line to number them within the section instead.
% \numberwithin{thm}{section}

% Convenience aliases so the AMS environment names read like the
% article-style names used elsewhere in this collection.
\newtheorem{theorem}{Theorem}
\newtheorem{lemma}[theorem]{Lemma}
\newtheorem{proposition}[theorem]{Proposition}
\newtheorem{corollary}[theorem]{Corollary}
\theoremstyle{definition}
\newtheorem{definition}[theorem]{Definition}
\newtheorem{example}[theorem]{Example}
\newtheorem{convention}[theorem]{Convention}
\theoremstyle{remark}
\newtheorem{remark}[theorem]{Remark}
\newtheorem{notation}[theorem]{Notation}

\allowdisplaybreaks[2]

% Project macros.
\newcommand{\N}{\mathbb{N}}
\newcommand{\Z}{\mathbb{Z}}
\newcommand{\Q}{\mathbb{Q}}
\newcommand{\R}{\mathbb{R}}
\newcommand{\C}{\mathbb{C}}
\newcommand{\F}{\mathbb{F}}
\DeclareMathOperator{\rank}{rank}
\DeclareMathOperator{\supp}{supp}
\DeclareMathOperator{\im}{im}
\DeclareMathOperator{\id}{id}
\DeclareMathOperator{\tr}{tr}
\DeclareMathOperator{\Spec}{Spec}
\DeclareMathOperator{\Hom}{Hom}
\DeclareMathOperator{\End}{End}
\newcommand{\dd}{\mathrm{d}}
\newcommand{\abs}[1]{\left\lvert #1 \right\rvert}
\newcommand{\norm}[1]{\left\lVert #1 \right\rVert}
\newcommand{\ip}[2]{\left\langle #1, #2 \right\rangle}


% Short forms in the optional argument drive the running head.
\title[Short running title]{A General Reusable AMS-Style Manuscript Template}

% \NoCaseChange protects bracketed or initial-only author names from the
% small-caps running head that amsart applies.
\author{\NoCaseChange{[OI]}}
\date{\today}

% In amsart, \subjclass and \keywords must be set BEFORE \maketitle:
% \maketitle emits them as the administrative footnotes of page 1.
% Setting them afterwards leaves them as loose body text.
\subjclass[2020]{Primary 00A00; Secondary 00A01}
\keywords{template, manuscript, AMS style}

\begin{document}

% In amsart the abstract is typeset BEFORE \maketitle.
\begin{abstract}
This template reproduces the AMS-flavoured manuscripts in the
collection: the \texttt{amsart} class with \texttt{reqno}, an abstract
placed before \texttt{\textbackslash maketitle}, AMS subject
classification and keywords, and \texttt{amsplain} bibliography style.
Replace this text with one self-contained paragraph stating the
problem, the main result, and its significance.
\end{abstract}

\maketitle

\tableofcontents
\clearpage

% ---------------------------------------------------------------------
% sections/introduction.tex
% ---------------------------------------------------------------------

\section{Introduction}
\label{sec:introduction}

State the problem, fix the notation, and state the main theorem.
Then place the result against the literature with citations such as
\cite{example-classical} and \cite{example-modern}, saying precisely
which hypothesis is relaxed.

\begin{theorem}[Main theorem]
\label{thm:main}
Let $X$ be a compact metric space and let $\mu$ be a Borel probability
measure on $X$. Then
\begin{equation}
\label{eq:main}
  \int_X f \dd\mu \le \norm{f}_{L^2(\mu)}
\end{equation}
for every $f \in L^2(\mu)$.
\end{theorem}

\begin{proof}
Apply Cauchy--Schwarz to the pairing of $f$ with the constant function
$1$, using $\mu(X) = 1$.
\end{proof}

\begin{remark}
\label{rem:scope}
The compactness hypothesis is used only to guarantee that
$L^2(\mu)$ is separable.
\end{remark}

The remainder of the paper is organised as follows.
\Cref{sec:geometry} proves the supporting estimate, and
\Cref{sec:conclusion} records consequences.

% ---------------------------------------------------------------------
% sections/geometry.tex
% ---------------------------------------------------------------------

\section{Supporting estimate}
\label{sec:geometry}

\begin{lemma}
\label{lem:aux}
Let $g \in L^\infty(\mu)$ with $\norm{g}_{L^\infty} \le 1$. Then
\begin{equation}
\label{eq:aux}
  \abs{\int_X g \dd\mu} \le 1 .
\end{equation}
\end{lemma}

\begin{proof}
Bound the integrand pointwise by $1$ and use $\mu(X) = 1$.
\end{proof}

\begin{proof}[Proof of \cref{thm:main}]
Approximate $f$ in $L^2(\mu)$ by bounded functions, apply
\cref{lem:aux} to the normalised approximants, and pass to the limit
using \cref{eq:aux}.
\end{proof}

% ---------------------------------------------------------------------
% sections/conclusion.tex
% ---------------------------------------------------------------------

\section{Consequences}
\label{sec:conclusion}

\begin{corollary}
\label{cor:main}
Under the hypotheses of \cref{thm:main}, the map
$f \mapsto \int_X f \dd\mu$ is continuous on $L^2(\mu)$.
\end{corollary}

\begin{proof}
Immediate from \cref{eq:main}.
\end{proof}

\subsection*{Open problems}

\begin{enumerate}
  \item Can compactness of $X$ be replaced by completeness?
  \item Is \cref{cor:main} sharp in the operator norm?
\end{enumerate}


\appendix
% ---------------------------------------------------------------------
% sections/appendix.tex
% \appendix in main.tex renumbers sections as A, B, ...
% ---------------------------------------------------------------------

\section{Deferred computation}
\label{app:deferred}

\begin{lemma}
\label{lem:app}
The constant in \cref{eq:main} cannot be improved.
\end{lemma}

\begin{proof}
Take $f = \mathbf{1}_X$; then both sides of \cref{eq:main} equal $1$.
\end{proof}


\bibliographystyle{amsplain}
\bibliography{references}

\end{document}
